
\documentclass{article}
\usepackage{colortbl}
\usepackage{makecell}
\usepackage{multirow}
\usepackage{supertabular}

\begin{document}

\newcounter{utterance}

\centering \large Interaction Transcript for game `taboo', experiment `medium\_en', episode 17 with reflexion{-}Qwen3.5{-}4B.
\vspace{24pt}

{ \footnotesize  \setcounter{utterance}{1}
\setlength{\tabcolsep}{0pt}
\begin{supertabular}{c@{$\;$}|p{.15\linewidth}@{}p{.15\linewidth}p{.15\linewidth}p{.15\linewidth}p{.15\linewidth}p{.15\linewidth}}
    \# & $\;$A & \multicolumn{4}{c}{Game Master} & $\;\:$B\\
    \hline

    \theutterance \stepcounter{utterance}  
    & & \multicolumn{4}{p{0.6\linewidth}}{
        \cellcolor[rgb]{0.9,0.9,0.9}{
            \makecell[{{p{\linewidth}}}]{
                \texttt{\tiny{[P1$\langle$GM]}}
                \texttt{You have attempted this same task before and did not succeed. These are the plans you wrote after those attempts.} \\
\texttt{Trial \#1: New plan: I failed because I provided clues that were too vague and allowed the opponent to guess related words like "scratch" instead of the target "nick". I should have used more specific descriptors that directly reference the target word's unique characteristics, such as mentioning it is a short, sharp mark often found on fingernails or teeth, rather than generic terms for surface damage. Next time, I will focus on distinguishing "nick" from "dent" and "cut" by emphasizing its small size and sharpness without using synonyms.} \\
\\ 
\texttt{Trial \#2: New plan: I failed because my clues were too descriptive of the *action* or *result* (like "bite" or "shallow mark") rather than the specific *object* itself, allowing the opponent to guess related words like "scratch" or "puncture." Next time, I will focus exclusively on the noun definition of "nick" as a specific type of small cut or notch, explicitly distinguishing it from "dent" (which implies depth) and "cut" (which implies length or severity), using phrases like "a small notch" or "a tiny cut on the surface" to force the opponent to identify the exact target word within the three attempts.} \\
\\ 
\texttt{Trial \#3: New plan: I failed because my clue described the *nature* of the mark (sharp, tiny) but failed to explicitly distinguish it from the related words "dent" and "cut" based on the provided definitions. The opponent likely guessed "cut" or "scratch" because my description was too generic. Next time, I will use a definition that strictly contrasts "nick" with "dent" by emphasizing that it is a **shallow** mark with **no depth**, and with "cut" by emphasizing it is a **small** or **minor** mark, explicitly stating it is not a deep dent or a long cut. I will phrase the clue as: "A shallow, small mark on a surface that is not deep enough to be a dent and not long enough to be a cut."} \\
\\ 
\\ 
\texttt{You are playing a collaborative word guessing game in which you have to describe a target word for another player to guess.} \\
\\ 
\texttt{Rules:} \\
\texttt{(a) You have to reply in the form: CLUE: <some text>. Guesses from the other player will start with GUESS.} \\
\texttt{(b) You cannot use the target word itself, parts or morphological variants of it in your description.} \\
\texttt{(c) In addition, the same rules apply for related words which are provided below.} \\
\\ 
\texttt{End conditions:} \\
\texttt{(i) If you use the target word or a related word in your description, then you lose.} \\
\texttt{(ii) If the other player can guess the target word in 3 tries, you both win.} \\
\\ 
\texttt{Let us start.} \\
\\ 
\texttt{This is the target word that you need to describe and that the other player needs to guess:} \\
\\ 
\texttt{nick} \\
\\ 
\texttt{Related words are:} \\
\\ 
\texttt{{-} groove} \\
\texttt{{-} dent} \\
\texttt{{-} cut} \\
\\ 
\texttt{Important: You are under time pressure, give short descriptions that are to the point!} \\
            }
        }
    }
    & \\ \\

    \theutterance \stepcounter{utterance}  
    & \multicolumn{4}{p{0.6\linewidth}}{
        \cellcolor[rgb]{0.9,0.9,0.9}{
            \makecell[{{p{\linewidth}}}]{
                \texttt{\tiny{[P1$\rangle$GM]}}
                \texttt{CLUE: A small, shallow mark on a surface that is not deep enough to be a dent and not long enough to be a cut.} \\
            }
        }
    }
    & & \\ \\

    \theutterance \stepcounter{utterance}  
    & & & \multicolumn{2}{p{0.3\linewidth}}{
        \cellcolor[rgb]{0.9,0.9,0.9}{
            \makecell[{{p{\linewidth}}}]{
                \texttt{\tiny{[GM$|$GM]}}
                \texttt{continue} \\
            }
        }
    }
    & & \\ \\

    \theutterance \stepcounter{utterance}  
    & & & \multicolumn{2}{p{0.3\linewidth}}{
        \cellcolor[rgb]{0.9,0.9,0.9}{
            \makecell[{{p{\linewidth}}}]{
                \texttt{\tiny{[GM$|$GM]}}
                \texttt{Related word 'dent' (stem=dent) is similar to clue word 'dent' (stem=dent)} \\
            }
        }
    }
    & & \\ \\

    \theutterance \stepcounter{utterance}  
    & \multicolumn{4}{p{0.6\linewidth}}{
        \cellcolor[rgb]{0.95,0.95,0.95}{
            \makecell[{{p{\linewidth}}}]{
                \texttt{\tiny{[P2$\rightarrow$P1]}}
                \texttt{Post{-}episode Reflexion (not sent to GM):} \\
\texttt{New plan: I failed because my previous clue was too generic and allowed the opponent to guess related words like "scratch" or "dent" due to a lack of precise boundary conditions. Next time, I will construct a clue that explicitly defines "nick" as a **shallow** mark with **no depth** (distinguishing it from "dent") and a **short** mark (distinguishing it from "cut"), while strictly avoiding synonyms like "scratch" or "puncture." I will phrase the clue as: "A shallow, short mark on a surface that has no depth and is not long enough to be considered a cut."} \\
            }
        }
    }
    & & \\ \\

\end{supertabular}
}

\end{document}
